\scanpage{34}
\mspage{15}
\[
(53)\qquad a_i^\omega\varphi_i^\omega(g)=\varphi_i^{\omega'}(\varepsilon_0(g))
\]
La donnée d'un endomorphisme du groupoïde \inserted{profini}, fixant les sommets et qui commute à $G_J$, équivaut à celle des \deleted{éléments} transformés des $a_i^\omega,c_i^\omega,d_i^\omega$, de façon \deleted{qui} compatible à l'action \unsure{de|des} $G_J$, et de façon à satisfaire aux relations (48).

On aura
\[
(54)\qquad\left\{\begin{aligned}
u(a_i^\omega)&=a_i^\omega\varphi_i^\omega(\lambda)\\
u(c_i^\omega)&=\varphi_{\omega'i}^{\omega'}(\eta)c_i^\omega=c_i^\omega\psi_{\unsure{0|i}}^\omega(\eta)\\
u(d_i^\omega)&=\psi_i^\omega(\xi)d_i^\omega=d_i^\omega\widetilde\psi_i^\omega(\xi)
\end{aligned}\right.
\]
avec $\lambda,\xi,\eta\in\widehat\Pi_{0,3}$, soumis à la seule relation déduite de (48). Au lieu de refaire le calcul, un peu fastidieux, pour expliciter la relation en $(\lambda,\eta,\xi)$, je préfère noter que le \deleted{le} groupoïde envisagé est la réunion des groupoïdes $\mathrm{II}$ et $\mathrm{IV}$, dont l'intersection est $\mathrm{III}$, donc \inserted{la donnée d'}un endom. \inserted{(fixant les sommets et comm. à $G_J$)} revient à celle d'un \inserted{tel} endom. $(\lambda,\eta)$ de $\mathrm{II}$, et d'un autre tel endom. \inserted{$\xi$} de $\mathrm{IV}$, de telle façon qu'ils se recollent \deleted{\unsure{en} \unsure{dans} $\mathrm{III}$} par un $=$ endomorphisme $(\widetilde\alpha)$ de $\mathrm{III}$, i.e tels qu'on ait à la fois
\[
(55)\qquad\widetilde\alpha=\eta^{-1}\varepsilon_0(\lambda\eta)\qquad\text{et}\qquad
\widetilde\alpha=\xi\rho^{-1}(\xi^{-1})
\]
\margin{33}
